% ============================================================================
% figures/gantt.tex
% Planning GANTT — remarque encadrante : « Vous devez ajouter la planif GANTT »
% Les jalons sont positionnés à partir de l'historique Git réel du dépôt
% (premier commit dataset : 10 février 2026 ; dernier commit rapport : 23 avril 2026).
% ============================================================================
\begin{figure}[H]
\centering
\begin{ganttchart}[
  hgrid, vgrid,
  x unit=0.52cm,
  y unit chart=0.55cm,
  y unit title=0.6cm,
  canvas/.style={draw=black!30, dash pattern=on 1pt off 1pt},
  bar/.style={fill=blue!40!black, draw=blue!60!black, rounded corners=1pt},
  bar label font=\footnotesize\sffamily,
  bar top shift=0.15,
  bar height=0.6,
  group/.style={fill=blue!70!black, draw=blue!80!black},
  group left shift=0,
  group right shift=0,
  group peaks width=0,
  group label font=\footnotesize\sffamily\bfseries,
  milestone/.style={fill=orange!70!black, draw=orange!80!black},
  milestone label font=\scriptsize\sffamily\itshape,
  title label font=\scriptsize\sffamily,
  title/.style={draw=black!30, fill=black!5},
  today=8,
  today rule/.style={draw=red!70, thick, dashed},
  today label=.,
  today label font=\tiny,
]{1}{11}

  % En-têtes : semaines
  \gantttitle{Février 2026}{2}
  \gantttitle{Mars 2026}{4}
  \gantttitle{Avril 2026}{5} \\
  \gantttitlelist{1,...,11}{1} \\

  % ---------- Phase 1 : Cadrage ----------
  \ganttgroup{Phase 1 --- Cadrage}{1}{1} \\
  \ganttbar{Étude de l'existant}{1}{1} \\
  \ganttbar{Spécification fonctionnelle}{1}{1} \\

  % ---------- Phase 2 : Dataset ----------
  \ganttgroup{Phase 2 --- Dataset}{1}{2} \\
  \ganttbar{Investisseurs (41)}{1}{2} \\
  \ganttbar{Subventions (25) \& métriques (18)}{1}{2} \\
  \ganttbar{FAQ (54), guides, régulations}{1}{2} \\
  \ganttmilestone{V3 dataset figé}{2} \\

  % ---------- Phase 3 : Pipeline RAG ----------
  \ganttgroup{Phase 3 --- Pipeline RAG}{2}{3} \\
  \ganttbar{\code{data\_loader.py}, \code{config.py}}{2}{2} \\
  \ganttbar{Embeddings + ingest + chunking}{2}{3} \\
  \ganttbar{ChromaDB + retriever}{3}{3} \\

  % ---------- Phase 4 : Modules métier ----------
  \ganttgroup{Phase 4 --- Modules métier}{3}{3} \\
  \ganttbar{Scoring, diagnostic}{3}{3} \\
  \ganttbar{Matching investisseurs \& subventions}{3}{3} \\
  \ganttbar{Routeur \& tests unitaires}{3}{3} \\

  % ---------- Phase 5 : Backend + Frontend ----------
  \ganttgroup{Phase 5 --- API \& interface}{3}{5} \\
  \ganttbar{FastAPI (endpoints sync)}{3}{3} \\
  \ganttbar{Frontend Next.js initial}{3}{4} \\
  \ganttbar{Sidebar + conversations}{4}{4} \\
  \ganttbar{SSE streaming + sources}{5}{5} \\
  \ganttbar{Dark mode, guided tour, responsive}{4}{5} \\
  \ganttmilestone{MVP fonctionnel}{5} \\

  % ---------- Phase 6 : Rapport ----------
  \ganttgroup{Phase 6 --- Rédaction \& soutenance}{6}{11} \\
  \ganttbar{Rédaction v1 (Markdown)}{6}{8} \\
  \ganttbar{Retours encadrante}{8}{9} \\
  \ganttbar{Refonte LaTeX + GANTT + figures}{9}{11} \\
  \ganttbar{Préparation soutenance}{10}{11} \\
  \ganttmilestone{Remise version finale}{11} \\

\end{ganttchart}
\caption{Planning prévisionnel et réel du projet Tamwil~AI, de février à avril 2026. Les phases 2 à 5 ont été réalisées sur cinq semaines (commits Git entre le 10 février et le 25 mars 2026), la phase 6 couvre la rédaction et l'itération avec l'encadrante jusqu'à la soutenance.}
\label{fig:gantt}
\end{figure}
